\scanpage{10}
Rectificatif sur critère d'irréductibilité géom.

\diagram{\begin{tikzcd}[ampersand replacement=\&,column sep=1.6em,row sep=.45em]
{} \&[-.3em] {} \&[-.5em] Z'\arrow[dr,dash,"{f'}"'] \& {}\\
Z\arrow[r,"f"] \& X\arrow[d,dash] \& {} \& X'\arrow[ll,"p"']\arrow[d,dash]\\
{} \& k\arrow[rr,dash] \& {} \& k'\\
{} \& {} \& {} \& \deleted{\text{\unsure{sép clos|séparable|sépar. clos}}}
\end{tikzcd}}
Supposons a) $Z$ géom. irréd, b) $f(Z)$ non contenu dans l'ens des points \deleted{\illegible{1}} où passent plusieurs comp. irréd. géom (p. ex. $f(Z)$ contient un pt géom normal, \deleted{\unsure{ou|un}} \inserted{voire} géom régulier, de $X_{\mathrm{red}}$). Alors $X$ géom irréductible.

\uline{Dém} [OPS $Z$ réduit à un pt \deleted{\unsure{de}} $x$ de $X$.] \deleted{\unsure{Alors $X$ irréd|On a $X$ irréd|Alors $X$ étant}} Soit $k'$ \deleted{clôture séparable} \inserted{ext finie \unsure{gal.|galois.}} de $k$, \deleted{\unsure{Alors|Soit}} $X'_i$ une composante irréductible de $X'$.  \deleted{\unsure{Elle domine|Elle devient}} contenant $f'(Z')$\inserted{(qui est \unsure{irréd|irrég})}, les \unsure{transformés|transformées} par $\pi$ sont des composantes irréductibles contenant $f'(Z')$, donc sont \uline{égales} à $\deleted{\unsure{f(|g(}}X'_i$.  Donc $X'_i$ est l'image inverse de une partie fermée irréductible de $X$, savoir de $\unsure{p|p'|f}(X'_i)$, or $p(X'_i)=X$ [platitude \deleted{\unsure{de|du}} \unsure{]|)} $p$ + $p$ fermé], donc $X'_i=X'$, donc $X'$ est irréd.

\unsure{Contre'exemple|Contr'exemple} si on laisse tomber b)

\[
X=\mathbb{R}[x,y]/(x^{\unsure{2|7}}+y^2),\quad Z=\operatorname{Spec}\mathbb{R},
\]
\diagram{\begin{tikzcd}[ampersand replacement=\&,row sep=.2em,column sep=1.5em]
Z\arrow[r] \& X\\
\mathbb{R} \& \mathbb{R}[x,y]/{-}\arrow[l]\\
0 \& x,y\arrow[l,maps to]
\end{tikzcd}}
\par\noindent\hfill\rule{.25\linewidth}{.4pt}\hfill\strut\par
\deleted{\unsure{Complément|Compléments} sur $\pi_0(X/k)$}

\uline{Remarque} Dans l'équivalence $\mathrm{Et}(k)\simeq(\pi\text{-}\mathrm{Ens})$, \deleted{les \unsure{ensembles}} \deleted{\unsure{on voit|on a vu}} \margin{\deleted{$\unsure{k'\underset{\pi}{\longleftarrow}k|k'\underset{\pi}{-}k'}$}}

\diagram{\begingroup\small\(\left\{\begin{array}{@{}l@{\;\;}l@{}}
\multicolumn{1}{c}{X} & \multicolumn{1}{c}{E}\\
\parbox[t]{.43\linewidth}{$X$ \unsure{quasi-compact|quasicompact} (de degré $n$)} & \parbox[t]{.44\linewidth}{--- $E$ fini (\unsure{dé|de|d'é} ordre $n$)}\\[.5ex]
\parbox[t]{.43\linewidth}{\deleted{Décomposition de $X$ en}} & {}\\
\parbox[t]{.43\linewidth}{$X$ connexe [ou : irréductible], \inserted{ou $X\simeq\operatorname{Spec}\unsure{K|k'|K'}$, $K$ ext. finie sép.\unsure{]|)}}} & \parbox[t]{.44\linewidth}{--- $\pi$ transitif sur $E$}\\[.6ex]
\parbox[t]{.43\linewidth}{Décomposition de $X$ en somme de \unsure{compos.|composantes} \deleted{\unsure{disjointes|disj.}} connexes $\ne$} & \parbox[t]{.44\linewidth}{--- décomposition de $E$ en sommes d'orbites}\\[.6ex]
\parbox[t]{.43\linewidth}{$\varprojlim$ finies de schémas} & \parbox[t]{.44\linewidth}{--- $\varprojlim$ finies de $\pi$-ens}\\
\parbox[t]{.43\linewidth}{$\varinjlim$ \deleted{\unsure{fin|finies}} de schémas.} & \parbox[t]{.44\linewidth}{--- $\varinjlim$ de $\pi$-ens.}\\
\parbox[t]{.43\linewidth}{\unsure{mono|momo} $X\to Y$ \inserted{(resp. épi)}} & \parbox[t]{.44\linewidth}{--- $\unsure{E-F|E\to F}$ injectif (resp. surjectif)}\\
\parbox[t]{.43\linewidth}{[$X$ un torseur \inserted{à dr.} sous $G$ \deleted{\unsure{groupe f|groupe}} groupe fini ordinaire} & \parbox[t]{.44\linewidth}{$E$ un torseur à droite sous $G$; les opérations de $G$ commutant à l'action de $\pi$ (donc si on choisit un \unsure{$e\in E$|$e$ de $E$}, la structure est définie par un hom continu $\pi\to G$)]}
\end{array}\right.\)\endgroup}
\uline{Rappel sur la construction de} $\pi_0(X/k)$ lorsque \deleted{\unsure{pour des extensions|pour les extensions}} \deleted{\unsure{pour|génér.}} $X_{k'}$ ($k'$ clôture parfaite de $k$) est localement connexe, \deleted{\unsure{donc|dans}} \inserted{et} ses \deleted{\unsure{$X_{k'}$|$X_k$}} composantes connexes définissables sur des \inserted{\unsure{ss.|sép.|sé.}} ext. finies de $k$ (p. ex. $X$ qu. cpct, ou $X$ loc. de \uline{\unsure{pr. finie|p. finie|t. f.|t. fini}} sur $k$)

\mspage{2}
% archive sheet number: 9
